Nuevas tecnologías emergentes hacen uso de la hibridación entre neuronas vivas y artificiales, formando circuitos que dependen de estas células especializadas para favorecer el estudio de su dinámica y de los fenómenos involucrados en su coordinación funcional. Distintos sistemas operativos en tiempo real y otras herramientas de modelado neuronal actuales pueden facilitar la interacción entre neuronas vivas y artificiales, por lo que se proponen algunas opciones en el contexto del estado del arte, así como pruebas de rendimiento y latencia para los mismos.

El objetivo de las interacciones deberá ser poder formar un ritmo funcional en el circuito neuronal biohíbrido, con una parte biológica (la neurona viva) y una parte artificial (el modelo neuronal), para su posterior uso en el ámbito de investigación. Para ello, se deben respetar las restricciones temporales que presentan dichas interacciones, de lo contrario, no se podrá establecer una conexión realista entre la parte biológica y la parte artificial. En dichas restricciones, es imperativo el uso de sistemas y herramientas especializados en este tipo de problemas basados en tiempo real.

En consecuencia, se presenta en primera instancia la formalización de circuitos biohíbridos para su uso en tiempo real, junto con una pequeña recapitulación del estado actual de las herramientas \textit{realtime}. Tras esto, se ahondará en el uso de una de las herramientas más utilizadas en la comunidad de investigación en neurociencia, \textit{RTXI}, un programa que permite realizar experimentos de electrofisiología en tiempo real. Dicho programa ha sido adaptado en este trabajo para su uso en el kernel \textit{Preempt-RT} de Linux, el oficial de tiempo real de esta familia de sistemas operativos. También se presentan una serie de validaciones en forma de pruebas de latencia y creación de circuitos biohíbridos, para comprobar su correcto funcionamiento en el contexto de neurociencia experimental.

Las pruebas de latencia y rendimiento se han realizado tanto en \textit{RTXI} como en \textit{cyclictest}, un programa de pruebas de la suite de tests \textit{rt-test} de la \textit{Linux Foundation}. En cuanto a los circuitos biohíbridos, se ha realizado en primer lugar una validación con una neurona electrónica (implementada en un circuito analógico) que presenta las mismas características temporales que una neurona viva. En segundo lugar, el sistema se ha validado con un registro electrofisiológico de una neurona viva que actúa de estímulo a un modelo neuronal que se ejecuta en el sistema desarrollado.


